% BOOK-FILE-ID: f-5d9116720103
\subsection{隣接4項型}\label{節：隣接4項型}
更新に使う項を増やして,
\[
  a_{n+3}+pa_{n+2}+qa_{n+1}+ra_n=0
\]
のような\bookterm{recurrence}{漸化式}を,\,$n\ge1$で考える.
初めの三項を自由に選ぶ三階の\bookterm{homogeneous}{同次}\bookterm{linear}{線形}\bookterm{recurrence}{漸化式}なので,解全体は3\bookterm{dimension}{次元}である(\sectionref*{節：隣接3項型}).
\sectionref*{節：隣接3項型}と同じく,\,$d_{n+1}=\gamma d_n$を目標に,
\bookterm{auxiliary}{補助数列}$d_n=a_{n+2}+\alpha a_{n+1}+\beta a_n$を作る.
その条件は
\BookMathLayout{\[
    a_{n+3}+\alpha a_{n+2}+\beta a_{n+1}=\gamma \left(a_{n+2}+\alpha a_{n+1}+\beta a_{n}\right)
  \]}{
  \[
    \begin{aligned}
      a_{n+3}+\alpha a_{n+2}+\beta a_{n+1}\\
      {}=\gamma\left(a_{n+2}+\alpha a_{n+1}+\beta a_n\right)
    \end{aligned}
  \]
}
である.両辺を左側に集め,元の\bookterm{recurrence}{漸化式}と係数を比べると,
\[
  \alpha-\gamma=p,\qquad
  \beta-\gamma\alpha=q,\qquad
  -\gamma\beta=r
\]
となる.最初の二式から
\[
  \alpha=p+\gamma,\qquad
  \beta=q+p\gamma+\gamma^2
\]
なので,最後の式へ代入すると
\[
  \gamma^3+p\gamma^2+q\gamma+r=0
\]
を得る.実数解を一つ選べば,\,$\alpha,\beta$も決まる.\footnote{三次方程式は少なくとも一つ実数解を持つが,計算しやすい値を見つけられるかは別の問題である.}
例題では因数分解や\bookterm{difference}{階差}を手掛かりに選ぶ.
$d_1=a_3+\alpha a_2+\beta a_1$から$d_n$を求めると,
\[
  a_{n+2}+\alpha a_{n+1}+\beta a_n=d_1\gamma^{\,n-1}
\]
という3項間\bookterm{recurrence}{漸化式}が得られる.\footnote{$\gamma=0$の場合は$d_1$を初項として,\,$n\ge2$では$d_n=0$と読む.
  また,\,\sectionref{節：定数の付加項がある場合}で扱ったのは右辺が定数$r$の場合であり,ここでは右辺が指数関数$d_1\gamma^{n-1}$になる点が異なる.}
$d_n$を求めた後,\,$a_n$へ戻す.
\[
  \begin{array}{c|c}
    \text{求める量}&\text{使う初期値}\\ \hline
    d_n&d_1\\
    a_n&a_1,a_2
  \end{array}
\]
$a_3$の情報は$d_1$に入っている.
得られた式は\sectionref{節：階差等比複合型}の考え方で,右辺と同じ更新をする量を引いて解ける場合がある.
例えば第1問では,\,$p=-3,q=3,r=-1$なので上の方程式は$(\gamma-1)^3=0$となり,
$\gamma=1,\alpha=-2,\beta=1$を得る.
こうして$d_n=a_{n+2}-2a_{n+1}+a_n$が決まる.
元の式を\mbox{三\bookterm{difference}{階差}}と見ても,同じ量が得られる.

\noindent \bookblocklink{example-four-term-1}{problem}{answer}{【例題】}\par\nobreak
\begin{enumerate}[label=(\arabic*),parsep=3pt,format=\bookitemlink{example-four-term-1}{problem}{answer}]
  \item $\displaystyle a_1=1,\,a_2=2,\,a_3=4$\\
    $a_{n+3}=3a_{n+2}-3a_{n+1}+a_n\quad(n\ge1)$
  \item $\displaystyle a_1=2,\,a_2=1,\,a_3=5$\\
    $a_{n+3}=2a_{n+2}+a_{n+1}-2a_n\quad(n\ge1)$
  \item $\displaystyle a_1=1,\,a_2=4,\,a_3=11$\\
    $a_{n+3}=5a_{n+2}-7a_{n+1}+3a_n\quad(n\ge1)$
\end{enumerate}

\noindent\bookblocklink{example-four-term-1}{hint}{problem}{【方針】}\par\nobreak
\begin{enumerate}[label=(\arabic*),parsep=3pt,format=\bookitemlink{example-four-term-1}{hint}{problem}]
  \item $d_n=a_{n+2}-2a_{n+1}+a_n$とおき,\bookterm{difference}{階差}をさらにとると一定になることを利用します.
  \item $d_n=a_{n+2}-a_{n+1}-2a_n$とおくと$d_{n+1}=d_n$となります.\bookterm{initial}{初期条件}から$d_n=0$を導き,\,3項間\bookterm{recurrence}{漸化式}として解きます.
  \item $d_n=a_{n+2}-2a_{n+1}+a_n$とおくと$d_{n+1}=3d_n$となります.得られた右辺を満たす\bookterm{sequence}{数列}を引き,\bookterm{difference}{階差}が一定の形に帰着させます.
\end{enumerate}

\noindent\bookblocklink{example-four-term-1}{answer}{problem}{【解答】}\par\nobreak
\begin{enumerate}[label=(\arabic*),parsep=3pt,format=\bookitemlink{example-four-term-1}{answer}{problem}]
  \item $\displaystyle a_n=\frac{n^2-n+2}{2}$
  \item $\displaystyle a_n=2^{n-1}+(-1)^{n-1}$
  \item $\displaystyle a_n=n-1+3^{n-1}$
\end{enumerate}

\noindent\bookblocklink{example-four-term-1}{solution}{problem}{【解法】}\par\nobreak
\begin{enumerate}[label=(\arabic*),parsep=3pt,format=\bookitemlink{example-four-term-1}{solution}{problem}]
  \item $d_n=a_{n+2}-2a_{n+1}+a_n$と置くと,
    \[
      d_{n+1}-d_n=a_{n+3}-3a_{n+2}+3a_{n+1}-a_n=0
    \]
    より$d_n=d_1=1$である.
    $e_n=a_{n+1}-a_n$と置けば,\,$e_{n+1}-e_n=1,e_1=1$より$e_n=n$.従って,
    \[
      a_n=a_1+\sum_{k=1}^{n-1}e_k
      =1+\sum_{k=1}^{n-1}k
      =\frac{n^2-n+2}{2}
    \]
    を得る.
  \item
    \[
      d_n=a_{n+2}-a_{n+1}-2a_n
    \]
    とおくと,
    \[
      d_{n+1}-d_n=a_{n+3}-2a_{n+2}-a_{n+1}+2a_n=0
    \]
    である.また,\,$d_1=5-1-2\cdot2=0$だから$d_n=0$であり,
    \[
      a_{n+2}-a_{n+1}-2a_n=0
    \]
    となる.\bookterm{characteristic}{特性方程式}$x^2-x-2=0$の解は$2,-1$なので,
    \[
      a_n=A\,2^{n-1}+B(-1)^{n-1}
    \]
    とおける.\bookterm{initial}{初期条件}より$A+B=2,\,2A-B=1$であるから$A=B=1$を得る.
  \item
    \[
      d_n=a_{n+2}-2a_{n+1}+a_n
    \]
    とおくと,与えられた\bookterm{recurrence}{漸化式}から$d_{n+1}=3d_n$である.\bookterm{initial}{初期条件}より$d_1=11-2\cdot4+1=4$なので,
    \[
      a_{n+2}-2a_{n+1}+a_n=4\cdot3^{n-1}
    \]
    となる.一方,\,$3^{n-1}$の2\bookterm{difference}{階差}も$4\cdot3^{n-1}$である.そこで$b_n=a_n-3^{n-1}$とおくと,
    \[
      b_{n+2}-2b_{n+1}+b_n=0
    \]
    となる.\,$b_1=0,\,b_2=1$より$b_n=n-1$であるから,
    \[
      a_n=b_n+3^{n-1}=n-1+3^{n-1}
    \]
    を得る.
\end{enumerate}
